\chapter{Testing}

\tikzset{pbox/.style={draw=cgray,fill=cgray!8,rounded corners=2pt,font=\sffamily\small,
    inner sep=4pt,align=center},
  parr/.style={-{Stealth[length=2mm]},thick,cgray}}

Even a tiny change can break code that used to work. A test is a small program that calls your code with
known input and checks the output, so a break shows up in seconds, not in production. This chapter covers
the tools that check code without running it (linters), the two standard-library test tools, \code{pytest}
(the one you will use), plus fakes for API calls, property tests and automation.

% ------------------------------------------------------------
\section{The ladder of checks}

Printing values and eyeballing them does not scale. Use layers, cheapest first:

\begin{center}
%%MD-BEGIN
% ```mermaid
% flowchart LR
%     A["write code"] --> B["linter<br/>ruff / pylint<br/>no run"]
%     B --> C["type checker<br/>mypy<br/>no run"]
%     C --> D["tests<br/>pytest / unittest<br/>runs your code"]
%     D --> E["CI<br/>all of it, every push"]
% ```
%%MD-END
\begin{tikzpicture}[node distance=5mm]
  \node[pbox,draw=cblue,fill=cblue!10] (a) {write\\code};
  \node[pbox,right=of a] (b) {\textbf{linter}\\{\footnotesize ruff, pylint}\\{\footnotesize \emph{does not run}}};
  \node[pbox,right=of b] (c) {\textbf{type checker}\\{\footnotesize mypy}\\{\footnotesize \emph{does not run}}};
  \node[pbox,right=of c,draw=cgold,fill=cgold!10] (d) {\textbf{tests}\\{\footnotesize pytest, unittest}\\{\footnotesize \emph{runs your code}}};
  \node[pbox,right=of d,draw=cgreen,fill=cgreen!10] (e) {\textbf{CI}\\{\footnotesize all checks,}\\{\footnotesize every push}};
  \foreach \x/\y in {a/b,b/c,c/d,d/e} \draw[parr] (\x)--(\y);
\end{tikzpicture}
\end{center}

\begin{center}
\begin{tabularx}{\linewidth}{@{}l >{\raggedright\arraybackslash}X >{\raggedright\arraybackslash}X@{}}
\toprule
\textbf{Layer} & \textbf{Finds} & \textbf{Misses} \\
\midrule
Linter & undefined names, unused imports, style problems & wrong logic \\
Type checker & wrong types passed or returned (Chapter~14) & wrong values \\
Tests & wrong behaviour on the inputs you thought of & inputs you did not think of \\
\bottomrule
\end{tabularx}
\end{center}

% ------------------------------------------------------------
\section{Linters: Pylint, Pyflakes and Ruff}

A \textbf{linter} reads source code and reports likely errors and style problems without running it.
Take this three-line file. It has a typo: \code{totl} instead of \code{total}.

\begin{pycode}
# style1.py
import os
total = 1
print(totl)
\end{pycode}

Running it would crash with \code{NameError} only when line 3 executes. A linter reports it at once.

\subsection{Ruff (real output)}

\textbf{Ruff} is a very fast linter and formatter written in Rust. Install it with \code{python -m pip install ruff}
(or \code{uv}). Real output, with its default rule set for this version:

\begin{rawcode}
> ruff check --output-format concise style1.py
style1.py:1:1: I001 [*] Import block is un-sorted or un-formatted
style1.py:1:8: F401 [*] `os` imported but unused
style1.py:3:7: F821 Undefined name `totl`
Found 3 errors.
[*] 2 fixable with `--fix`.
\end{rawcode}

Read it as \code{file:line:column: CODE message}. \code{F821} is the real bug. The \code{[*]} marks problems Ruff can fix
itself:

\begin{rawcode}
> ruff check --fix style1.py           (removes the unused import, keeps the bug)
\end{rawcode}

Ruff also formats code, like a strict spell-checker for whitespace. Given an untidy file:

\begin{pycode}
def add(a,b):
    return a+b
x=add( 1,2 )
\end{pycode}

\code{ruff format messy.py} rewrites it to:

\begin{pycode}
def add(a, b):
    return a + b


x = add(1, 2)
\end{pycode}

\subsection{Pylint, Pyflakes, pycodestyle}

\begin{center}
\begin{tabularx}{\linewidth}{@{}l >{\raggedright\arraybackslash}X >{\raggedright\arraybackslash}X@{}}
\toprule
\textbf{Tool} & \textbf{Checks} & \textbf{Character} \\
\midrule
\code{pyflakes} & errors only: undefined names, unused imports & quick, few opinions \\
\code{pycodestyle} (was \code{pep8}) & layout against PEP~8: spacing, blank lines & style only \\
\code{pylint} & errors, style, naming, design smells, and a score out of 10 & thorough and noisy; configure it \\
\code{ruff} & most of the three above plus a formatter and import sorting & fast; the modern default \\
\bottomrule
\end{tabularx}
\end{center}

Pylint prefixes each message with a letter: \code{E} error, \code{W} warning, \code{C} convention (style),
\code{R} refactor, \code{F} fatal. For the file above its report has the shape below
(\emph{typical form; Pylint is not installed here, so not run}). An \code{E} line is the one to fix first.

\begin{rawcode}
************* Module style1
style1.py:1:0: C0114: Missing module docstring (missing-module-docstring)
style1.py:3:6: E0602: Undefined variable 'totl' (undefined-variable)
\end{rawcode}

\begin{gamerule}
Fix errors (\code{E}, \code{F}) first, style later. Do not chase a perfect score. Run \code{ruff check} and
\code{ruff format} before every commit; your editor can do it on save.
\end{gamerule}

% ------------------------------------------------------------
\section{Unittest}

Linters cannot know what your function \emph{should} return. For that, write a test. \code{unittest} is in the
standard library and needs no install.

\subsection{Arrange, act, assert}

Every test has three steps. The \emph{assertion} is the claim that must be true; if it is not, the test fails.

\begin{center}
%%MD-BEGIN
% ```mermaid
% flowchart LR
%     A["ARRANGE<br/>text = '  Claude   Haiku '"] --> B["ACT<br/>result = normalize_model_name(text)"]
%     B --> C["ASSERT<br/>result == 'claude-haiku'"]
% ```
%%MD-END
\begin{tikzpicture}[node distance=9mm]
  \node[pbox,draw=cblue,fill=cblue!10,text width=4cm] (a) {\textbf{ARRANGE}\\{\ttfamily\footnotesize text = "  Claude   Haiku "}};
  \node[pbox,right=of a,draw=cgold,fill=cgold!10,text width=4.3cm] (b) {\textbf{ACT}\\{\ttfamily\footnotesize result = normalize\_model\_name(text)}};
  \node[pbox,right=of b,draw=cgreen,fill=cgreen!10,text width=3.4cm] (c) {\textbf{ASSERT}\\{\ttfamily\footnotesize result == "claude-haiku"}};
  \draw[parr] (a)--(b); \draw[parr] (b)--(c);
\end{tikzpicture}
\end{center}

Here is the function under test, a first and naive version. It makes a model name safe for file names.

\begin{pycode}
# names.py
def normalize_model_name(name):
    return name.lower().replace(" ", "-")
\end{pycode}

A test class derives from \code{unittest.TestCase}. Each method whose name starts with \code{test} is a test.

\begin{pycode}
# test_names.py
import unittest

import names


class TestNormalize(unittest.TestCase):
    def setUp(self):
        self.fn = names.normalize_model_name   # runs before EACH test

    def tearDown(self):
        pass                                   # runs after EACH test

    def test_one_word(self):
        self.assertEqual(self.fn("Haiku"), "haiku")

    def test_two_words(self):
        self.assertEqual(self.fn("Claude Haiku"), "claude-haiku")

    def test_extra_spaces(self):
        self.assertEqual(self.fn("  Claude   Haiku "), "claude-haiku")

    def test_rejects_none(self):
        with self.assertRaises(AttributeError):
            self.fn(None)


if __name__ == "__main__":
    unittest.main()
\end{pycode}

\begin{rawcode}
> python test_names.py
F...
======================================================================
FAIL: test_extra_spaces (__main__.TestNormalize.test_extra_spaces)
----------------------------------------------------------------------
Traceback (most recent call last):
  File "...\test_names.py", line 20, in test_extra_spaces
    self.assertEqual(self.fn("  Claude   Haiku "), "claude-haiku")
AssertionError: '--claude---haiku-' != 'claude-haiku'
- --claude---haiku-
? --       --     -
+ claude-haiku

----------------------------------------------------------------------
Ran 4 tests in 0.001s

FAILED (failures=1)
\end{rawcode}

Read the first line: one character per test, \code{.} passed and \code{F} failed. The \code{-} and \code{+} lines show
actual and expected, and the \code{?} line points to the differences. The test found a real bug.
Fix the function, then rerun:

\begin{pycode}
# names.py, fixed
def normalize_model_name(name):
    return "-".join(name.lower().split())    # split() drops all extra spaces
\end{pycode}
\begin{rawcode}
> python test_names.py
....
----------------------------------------------------------------------
Ran 4 tests in 0.000s

OK
\end{rawcode}

\begin{center}
%%MD-BEGIN
% ```mermaid
% flowchart LR
%     A["add a test<br/>for a case"] --> B["run: FAIL (red)"]
%     B --> C["fix the code"]
%     C --> D["run: PASS (green)"]
%     D --> A
% ```
%%MD-END
\begin{tikzpicture}[node distance=9mm]
  \node[pbox] (a) {add a test\\for one case};
  \node[pbox,right=of a,draw=cred,fill=cred!10] (b) {run:\\\textbf{FAIL}};
  \node[pbox,right=of b] (c) {fix the\\code};
  \node[pbox,right=of c,draw=cgreen,fill=cgreen!10] (d) {run:\\\textbf{PASS}};
  \draw[parr] (a)--(b); \draw[parr] (b)--(c); \draw[parr] (c)--(d);
  \draw[parr] (d.north) to[bend right=35] node[above,font=\scriptsize\sffamily]{next case} (a.north);
\end{tikzpicture}
\end{center}

\begin{typecheck}
A test class is a \emph{class} (Chapter~11) inheriting \code{TestCase}; \code{self.assertEqual(a, b)} is an inherited method.
It prints the difference when \code{a != b}. Plain \code{assert a == b} also works in unittest, but its failure message is poor.
\end{typecheck}

\begin{center}
\begin{tabular}{@{}l l@{}}
\toprule
\textbf{Method} & \textbf{Passes when} \\
\midrule
\code{assertEqual(a, b)} & \code{a == b} \\
\code{assertTrue(x)} / \code{assertFalse(x)} & \code{x} is truthy / falsy \\
\code{assertIn(a, seq)} & \code{a in seq} \\
\code{assertIsNone(x)} & \code{x is None} \\
\code{assertAlmostEqual(a, b)} & floats equal to 7 decimal places \\
\code{assertRaises(Err)} & the block raises \code{Err} (use \code{with}) \\
\bottomrule
\end{tabular}
\end{center}

Run all tests in a folder with \code{python -m unittest} (it finds files named \code{test*.py}); add \code{-v} for one line per test.

\begin{center}
%%MD-BEGIN
% ```text
% test_a:  setUp -> test_a -> tearDown
% test_b:  setUp -> test_b -> tearDown      (fresh start for every test)
% ```
%%MD-END
\begin{tikzpicture}[node distance=3mm]
  \tikzset{st/.style={pbox,font=\ttfamily\footnotesize,inner sep=3pt}}
  \node[st,draw=cblue,fill=cblue!10] (s1) {setUp};
  \node[st,right=of s1,draw=cgold,fill=cgold!10] (t1) {test\_a};
  \node[st,right=of t1,draw=cblue,fill=cblue!10] (d1) {tearDown};
  \node[st,right=8mm of d1,draw=cblue,fill=cblue!10] (s2) {setUp};
  \node[st,right=of s2,draw=cgold,fill=cgold!10] (t2) {test\_b};
  \node[st,right=of t2,draw=cblue,fill=cblue!10] (d2) {tearDown};
  \foreach \x/\y in {s1/t1,t1/d1,s2/t2,t2/d2} \draw[parr] (\x)--(\y);
\end{tikzpicture}
\end{center}

% ------------------------------------------------------------
\section{Doctest}

\code{doctest} runs the examples you put \emph{inside a docstring}. They look like a session at the
\code{>{}>{}>} prompt. The test is also documentation that cannot go out of date.

\begin{pycode}
# names2.py
def normalize_model_name(name):
    """Make a model name safe for file names.

    >>> normalize_model_name("Haiku")
    'haiku'
    >>> normalize_model_name("  Claude   Haiku ")
    'claude-haiku'
    >>> normalize_model_name("")
    ''
    """
    return "-".join(name.lower().split())


if __name__ == "__main__":
    import doctest
    doctest.testmod()
\end{pycode}

Run it. No output means every example passed. Add \code{-v} to see them all:

\begin{rawcode}
> python names2.py
> python names2.py -v
Trying:
    normalize_model_name("Haiku")
Expecting:
    'haiku'
ok
...
3 tests in 2 items.
3 passed and 0 failed.
Test passed.
\end{rawcode}

(Middle lines omitted.) You can also run it without the \code{if} block: \code{python -m doctest names2.py}.
A wrong expectation shows what Python really returned:

\begin{rawcode}
Failed example:
    twice(2)
Expected:
    5
Got:
    4
***Test Failed*** 1 failures.
\end{rawcode}

\begin{trap}
Doctest compares the \emph{printed text} exactly. The result \code{'haiku'} must be written with the quotes, and a float
such as \code{0.30000000000000004} must match digit for digit. Use doctest for short examples; use pytest for real test suites.
\end{trap}

% ------------------------------------------------------------
\section{Pytest}

\textbf{pytest} is the most popular test tool. Install it with \code{python -m pip install pytest}. Tests are plain
functions with plain \code{assert}; there is no class and no \code{self.assertEqual}.

\begin{rawcode}
> python -m pytest -q
\end{rawcode}

\begin{center}
\begin{tabularx}{\linewidth}{@{}l X@{}}
\toprule
\textbf{Feature} & \textbf{What you get} \\
\midrule
Discovery & finds \code{test\_*.py} files and \code{test\_*} functions below the current folder \\
Rich failures & a failing \code{assert} shows both sides and a diff, no special method needed \\
Fixtures & functions that prepare test data or resources, injected by name \\
Parametrization & run one test over a table of inputs \\
Built-ins & \code{tmp\_path}, \code{monkeypatch}, \code{capsys}, \code{pytest.raises}, \code{pytest.approx} \\
\bottomrule
\end{tabularx}
\end{center}

\subsection{A first test file}

\begin{pycode}
# costs.py
RATES = {"small": (0.10, 0.50), "large": (4.0, 20.0)}   # USD per 1M tokens

def token_cost(model, tokens_in, tokens_out=0):
    if tokens_in < 0 or tokens_out < 0:
        raise ValueError("token counts must be >= 0")
    rate_in, rate_out = RATES[model]
    return (tokens_in * rate_in + tokens_out * rate_out) / 1e6
\end{pycode}

\begin{pycode}
# test_costs_pytest.py   (needs: pytest)
import pytest

from costs import token_cost


def test_zero_tokens_cost_nothing():
    assert token_cost("small", 0, 0) == 0


def test_known_value():
    assert token_cost("small", 1_000_000, 0) == pytest.approx(0.10)


def test_negative_tokens_rejected():
    with pytest.raises(ValueError, match="must be >= 0"):
        token_cost("small", -1)


def test_unknown_model():
    with pytest.raises(KeyError):
        token_cost("huge", 1)
\end{pycode}

Typical pytest report for a passing run (\emph{pytest is not installed in this environment, so this is the usual
format, not a captured run}; the file itself was syntax-checked and its logic verified by calling the functions):

\begin{rawcode}
============================= test session starts =============================
collected 4 items

test_costs_pytest.py ....                                               [100%]

============================== 4 passed in 0.03s ==============================
\end{rawcode}

On a failure, pytest prints the test source, a \code{>} on the failing line, and \code{E} lines with the values.

\begin{trap}
\textbf{Never compare floats with \code{==}.} Binary floats are approximate. Compare with \code{pytest.approx(x)}
(or \code{math.isclose} in plain code).
\begin{pycode}
import math
print(0.1 + 0.2 == 0.3, math.isclose(0.1 + 0.2, 0.3))
\end{pycode}
\begin{rawcode}
False True
\end{rawcode}
\end{trap}

\subsection{Parametrization: one test, many cases}

\code{parametrize} takes a string of argument names and a list of tuples; pytest runs the test once per tuple and reports
each separately.

\begin{pycode}
import pytest
from costs import token_cost

@pytest.mark.parametrize(
    "model, tin, tout, expected",
    [
        ("small", 1_000_000, 0, 0.10),
        ("small", 0, 1_000_000, 0.50),
        ("large", 1_000, 500, 0.014),
    ],
)
def test_token_cost_table(model, tin, tout, expected):
    assert token_cost(model, tin, tout) == pytest.approx(expected)
\end{pycode}

\subsection{Fixtures}

A \textbf{fixture} is a function marked \code{@pytest.fixture}. A test that lists its name as a parameter receives
its return value. Fixtures replace \code{setUp} and \code{tearDown}.

\begin{pycode}
# test_fixtures.py   (needs: pytest)
import sqlite3

import pytest


@pytest.fixture                      # scope="function" is the default
def history():
    return [{"role": "user", "content": "hi"}]


@pytest.fixture(scope="module")      # created once for this file
def calls_db():
    conn = sqlite3.connect(":memory:")
    conn.execute("create table calls (model text, cost real)")
    yield conn                       # the test runs here
    conn.close()                     # teardown, after the last test
\end{pycode}
The tests then list the fixtures they need by name:
\begin{pycode}
def test_insert(calls_db, history):
    calls_db.execute("insert into calls values ('small', 0.0001)")
    assert calls_db.execute("select count(*) from calls").fetchone()[0] == 1


def test_same_connection_reused(calls_db):
    assert calls_db.execute("select count(*) from calls").fetchone()[0] == 1


def test_save_reply(tmp_path):
    out = tmp_path / "reply.txt"     # tmp_path: a fresh temp folder
    out.write_text("hello")
    assert out.read_text() == "hello"
\end{pycode}

\begin{center}
%%MD-BEGIN
% ```text
% scope="function" (default):   test_a [fixture A1]   test_b [fixture A2]   test_c [fixture A3]
% scope="module":               test_a [---------- one fixture, shared ----------] test_c
% ```
%%MD-END
\begin{tikzpicture}[node distance=0mm]
  \tikzset{tt/.style={pbox,font=\ttfamily\footnotesize,minimum width=18mm,inner sep=3pt}}
  \node[font=\sffamily\footnotesize,anchor=east] at (-1.1,0.9) {\texttt{function}};
  \foreach \i/\nm in {0/test\_a,1/test\_b,2/test\_c}{
    \node[tt] (f\i) at (\i*4.6,0.9) {\nm};
    \node[pbox,draw=cgold,fill=cgold!10,font=\sffamily\scriptsize,minimum width=18mm] at (\i*4.6+2.2,0.9) {new fixture};}
  \node[font=\sffamily\footnotesize,anchor=east] at (-1.1,-0.3) {\texttt{module}};
  \node[pbox,draw=cgreen,fill=cgreen!10,font=\sffamily\scriptsize,minimum width=11cm,anchor=west] (big) at (-0.9,-0.3) {one fixture, created once, shared by all three tests};
  \foreach \i/\nm in {0/test\_a,1/test\_b,2/test\_c}{
    \node[tt] at (\i*4.6,-1.2) {\nm};
    \draw[parr,cgreen] (\i*4.6,-0.95)--(\i*4.6,-0.55);}
\end{tikzpicture}
\end{center}

\begin{gamerule}
Keep the default \code{function} scope so each test starts clean. Raise the scope (\code{module}, \code{session}) only for something
slow to build, like a database connection, and only if tests do not change it in ways that leak between them.
\end{gamerule}

\subsection{Handy flags}

\begin{center}
\begin{tabularx}{\linewidth}{@{}l X@{}}
\toprule
\textbf{Command} & \textbf{Effect} \\
\midrule
\code{python -m pytest -q} & quiet: one dot per test \\
\code{python -m pytest -v} & one line per test, with names \\
\code{python -m pytest -x} & stop at the first failure \\
\code{python -m pytest -k normalize} & run only tests whose name contains \emph{normalize} \\
\code{pytest test\_costs\_pytest.py::test\_unknown\_model} & run one test (start with \code{python -m}) \\
\code{python -m pytest -s} & show \code{print} output (otherwise captured) \\
\bottomrule
\end{tabularx}
\end{center}

\subsection{Testing code that calls the API}

A test must be fast, free and repeatable. A real API call is slow, costs money and gives a different answer
each time. So \textbf{pass the client in} as a parameter and give the test a fake one.

\begin{center}
%%MD-BEGIN
% ```mermaid
% flowchart LR
%     subgraph production
%       P1["summarize(client, text)"] --> P2["real Anthropic client"] --> P3["API (slow, costs money)"]
%     end
%     subgraph test
%       T1["summarize(client, text)"] --> T2["fake client (Mock)"] --> T3["canned reply"]
%     end
% ```
%%MD-END
\begin{tikzpicture}[node distance=6mm]
  \node[pbox,draw=cblue,fill=cblue!10,text width=2.8cm] (p1) {\ttfamily\footnotesize summarize(client, text)};
  \node[pbox,right=of p1,draw=cred,fill=cred!10] (p2) {real client};
  \node[pbox,right=of p2,draw=cred,fill=cred!10] (p3) {API\\{\footnotesize slow, costs money}};
  \node[pbox,draw=cblue,fill=cblue!10,text width=2.8cm,below=10mm of p1] (t1) {\ttfamily\footnotesize summarize(client, text)};
  \node[pbox,right=of t1,draw=cgreen,fill=cgreen!10] (t2) {fake client};
  \node[pbox,right=of t2,draw=cgreen,fill=cgreen!10] (t3) {canned reply\\{\footnotesize instant, free}};
  \draw[parr] (p1)--(p2); \draw[parr] (p2)--(p3); \draw[parr] (t1)--(t2); \draw[parr] (t2)--(t3);
  \node[font=\sffamily\scriptsize,text=cgray,anchor=south west] at (p1.north west) {in production};
  \node[font=\sffamily\scriptsize,text=cgray,anchor=south west] at (t1.north west) {in a test};
\end{tikzpicture}
\end{center}

\begin{pycode}
# summarize.py
def summarize(client, text):
    """Ask the model for a one-line summary. Returns '' if no text block."""
    response = client.messages.create(
        model="small-model",
        max_tokens=200,
        messages=[{"role": "user", "content": f"Summarize in one line: {text}"}],
    )
    for block in response.content:
        if block.type == "text":
            return block.text.strip()
    return ""
\end{pycode}

\code{unittest.mock.Mock} (standard library) creates an object that accepts any attribute or call. We set what
\code{create} returns and later check how it was called.

\begin{pycode}
# test_summarize.py
import unittest
from types import SimpleNamespace
from unittest.mock import Mock

from summarize import summarize


def fake_client(*blocks):
    client = Mock()
    client.messages.create.return_value = SimpleNamespace(content=list(blocks))
    return client
\end{pycode}
Each test builds a fake client with the blocks the ``model'' should return:
\begin{pycode}
class TestSummarize(unittest.TestCase):
    def test_returns_text(self):
        client = fake_client(SimpleNamespace(type="text", text=" Short. "))
        self.assertEqual(summarize(client, "long text"), "Short.")

    def test_skips_thinking_block(self):
        client = fake_client(SimpleNamespace(type="thinking", thinking="..."),
                             SimpleNamespace(type="text", text="Answer"))
        self.assertEqual(summarize(client, "x"), "Answer")

    def test_no_text_block(self):
        client = fake_client(SimpleNamespace(type="thinking", thinking="..."))
        self.assertEqual(summarize(client, "x"), "")

    def test_sends_one_request(self):
        client = fake_client(SimpleNamespace(type="text", text="ok"))
        summarize(client, "hello")
        client.messages.create.assert_called_once()
        kwargs = client.messages.create.call_args.kwargs
        self.assertEqual(kwargs["max_tokens"], 200)
        self.assertIn("hello", kwargs["messages"][0]["content"])


if __name__ == "__main__":
    unittest.main(verbosity=2)
\end{pycode}
\begin{rawcode}
test_no_text_block (__main__.TestSummarize.test_no_text_block) ... ok
test_returns_text (__main__.TestSummarize.test_returns_text) ... ok
test_sends_one_request (__main__.TestSummarize.test_sends_one_request) ... ok
test_skips_thinking_block (__main__.TestSummarize.test_skips_thinking_block) ... ok

----------------------------------------------------------------------
Ran 4 tests in 0.001s

OK
\end{rawcode}

Note that the same fake lets us test the ``thinking block first'' case without paying for a real model.
In pytest, the \code{monkeypatch} fixture does a similar job for environment variables, for example
\code{monkeypatch.setenv(...)} with a dummy value for \code{ANTHROPIC\_API\_KEY}.

\begin{trap}
A test that passes only when the network is up, or when today's date is right, or when the model answers a
certain way, is a \emph{flaky} test. Fake the clock, the network and the model. Never put a real API key in a test or commit it.
\end{trap}

% ------------------------------------------------------------
\section{Hypothesis: property-based tests}

Instead of listing inputs, describe the \emph{kind} of input and a rule that must always hold; \textbf{Hypothesis} generates many
cases (100 by default), including odd ones you would not write, and shrinks a failure to a minimal example.

\begin{pycode}
# test_props.py   (needs: hypothesis)
from hypothesis import given
from hypothesis import strategies as st

from costs import token_cost

counts = st.integers(min_value=0, max_value=10**7)


@given(counts, counts, counts)
def test_more_output_never_cheaper(tokens_in, tokens_out, extra):
    cheaper = token_cost("small", tokens_in, tokens_out)
    dearer = token_cost("small", tokens_in, tokens_out + extra)
    assert dearer >= cheaper
\end{pycode}

(Install with \code{python -m pip install hypothesis}; not installed here. I checked the rule itself against
100\,000 random inputs with plain \code{random}: it held.) If you let \code{counts} include negative numbers,
Hypothesis would report a counterexample at once, since \code{token\_cost} raises \code{ValueError} for them.
That is the point: it finds the cases you forgot.

% ------------------------------------------------------------
\section{Nox}

\textbf{Nox} runs your checks in fresh environments, for several Python versions or tool combinations. You describe the
sessions in a Python file named \code{noxfile.py} and run \code{nox}.

\begin{pycode}
# noxfile.py   (needs: nox)
import nox


@nox.session(python=["3.11", "3.12"])
def tests(session):
    session.install("-r", "requirements.txt", "pytest")
    session.run("pytest", "-q")


@nox.session
def lint(session):
    session.install("ruff")
    session.run("ruff", "check", ".")
\end{pycode}

Nox is the successor of \code{tox}. Use it when your package must work on more than one Python version; for a personal
project, a plain \code{pytest} is enough.

% ------------------------------------------------------------
\section{Continuous integration}

\textbf{CI} runs your checks automatically on a server each time code is pushed, so a break is caught
before it is merged.

\begin{center}
%%MD-BEGIN
% ```mermaid
% flowchart LR
%     A["git push"] --> B["CI server<br/>fresh machine"]
%     B --> C["install dependencies"] --> D["ruff + pytest"]
%     D --> E{"pass?"}
%     E -- yes --> F["green check<br/>safe to merge"]
%     E -- no --> G["red cross<br/>fix first"]
% ```
%%MD-END
\begin{tikzpicture}[node distance=6mm]
  \node[pbox,draw=cblue,fill=cblue!10] (a) {\ttfamily git push};
  \node[pbox,right=of a] (b) {CI server\\{\footnotesize fresh machine}};
  \node[pbox,right=of b] (c) {install\\dependencies};
  \node[pbox,right=of c] (d) {\ttfamily ruff, pytest};
  \node[pbox,right=of d,draw=cgreen,fill=cgreen!10,yshift=6mm] (e) {green check\\{\footnotesize safe to merge}};
  \node[pbox,right=of d,draw=cred,fill=cred!10,yshift=-6mm] (f) {red cross\\{\footnotesize fix first}};
  \foreach \x/\y in {a/b,b/c,c/d} \draw[parr] (\x)--(\y);
  \draw[parr] (d)--(e.west); \draw[parr] (d)--(f.west);
\end{tikzpicture}
\end{center}

A minimal GitHub Actions workflow (file \code{.github/workflows/tests.yml}), not run here:

\begin{rawcode}
name: tests
on: [push, pull_request]
jobs:
  test:
    runs-on: ubuntu-latest
    steps:
      - uses: actions/checkout@v4
      - uses: actions/setup-python@v5
        with:
          python-version: "3.12"
      - run: python -m pip install -r requirements.txt pytest ruff
      - run: ruff check .
      - run: python -m pytest -q
\end{rawcode}

\begin{center}
\begin{tabularx}{\linewidth}{@{}l X@{}}
\toprule
\textbf{System} & \textbf{Note} \\
\midrule
GitHub Actions & built into GitHub; free for public repositories; the easy start \\
Jenkins & self-hosted, written in Java; common in companies \\
Buildbot & self-hosted, written in Python \\
Travis CI, CircleCI & hosted services; free tiers for open source \\
\bottomrule
\end{tabularx}
\end{center}

\begin{gamerule}[: CI and secrets]
CI machines are not yours. Store an API key, if a job truly needs one, in the host's \emph{secrets} setting, never in the repository.
Most tests should not need a key: use fakes.
\end{gamerule}

\begin{boss}[: testing what you do not control]
The hard part of testing is not \code{assert}; it is \emph{isolation}. Code that reads the clock, the environment, a random
number or a web API gives a different answer on each run or costs money. The fix is one idea: do not reach out from inside
the function; \textbf{pass the dependency in} (the client, the current time, a random seed) and give the test a fake or a fixed value.
Then the test depends only on its inputs, which is what makes it fast, free and repeatable. Practise it on \code{summarize(client, text)}:
the function never creates its own client.
\end{boss}

\begin{practice}
\begin{enumerate}[leftmargin=*,itemsep=2pt]
\item Run \code{ruff check} on a file that uses a variable before defining it. What is the code of the error, and does the linter run your program?
\item Write three \code{unittest} tests for a function \code{clip(text, n)} that cuts text to at most \code{n} characters. Which edge case would you add first?
\item What do \code{setUp} and \code{tearDown} do, and which pytest feature replaces them?
\item Rewrite your three tests as one \code{parametrize} test. What do you gain?
\item Why is \code{assert cost == 0.014} a bad test, and what do you write instead?
\item Your function calls \code{client.messages.create}. Describe how you test it without internet or an API key.
\end{enumerate}
\end{practice}

\begin{achievement}{Green Build}
\begin{itemize}[leftmargin=*,itemsep=1pt]
\item Lint and format code with Ruff, and read an error code and its location.
\item Write unit tests in \code{unittest} and \code{pytest}, including edge cases and exceptions.
\item Use fixtures and \code{parametrize} to remove repeated test code.
\item Test API-calling code with a fake client, with no network and no key.
\item Describe how CI runs the same checks on every push.
\end{itemize}
\end{achievement}

%%MD-BEGIN
% **Related:** [[10 - Functions]] (what you test), [[11 - Objects]] (`TestCase` is a class), [[13 - Development Environment]] (install pytest, ruff), [[14 - Type Hints and Documentation]] (mypy, docstrings), [[16 - Debugging]] (when a test fails), [[23 - Data in a Box - Persistent Storage]] (sqlite fixture).
%%MD-END
